\documentclass[11pt]{article}
\usepackage{explorer}
\newcommand{\ExplorerLanguage}{EN}
\newcommand{\ExplorerNumber}{04}
\newcommand{\ExplorerKind}{Lesson and method}
\begin{document}
\ExplorerTitle{Advanced high school · extension}{Bézout and integer equations}
\ExplorerSection{Before you begin}
Euclidean algorithm and divisibility in \(\mathbb{Z}\).\par
\ExplorerSection{Ideas and method}
For positive integers a and b, the Euclidean algorithm computes \(d = gcd(a,b)\). Back-substitution expresses d as au + bv for integers u and v: a Bézout identity.\par
The equation ax + by = c has integer solutions exactly when d divides c. Necessity follows because d divides ax + by; sufficiency follows by multiplying a Bézout identity by \(\frac{c}{d}\).\par
If \((x_{0},y_{0})\) is one solution, all solutions are \(x = x_{0} + (\frac{b}{d})t\) and \(y = y_{0} - (\frac{a}{d})t\) for integer t. Conditions \(x \geq  0\) and \(y \geq  0\) must then be imposed separately.\par
\ExplorerSection{Worked example}
For \(15x + 21y = 6\), use \(21 = 15 + 6\) and \(15 = 2 \times  6 + 3\). Thus \(3 = 3 \times  15 - 2 \times  21\). One solution is \((6,-4)\), and all solutions are \(x = 6 + 7t, y = -4 - 5t\).\par
\ExplorerSection{Explore}
Try several values of t in the example and verify the equation exactly.\par
\end{document}
